\documentclass{article}

\usepackage[utf8]{inputenc}
\usepackage[T1]{fontenc}
\usepackage[german]{babel}
\usepackage{arxiv}
\usepackage{graphicx}
\usepackage{hyperref}

\title{Topologische Stabilität der Riemannschen Nullstellen im oktonischen Primzahlmodell des 60-Flächners}

\author{
	\textbf{Autor Name} \\
	Bamberg, Deutschland \\
	\texttt{email@beispiel.de} \\
}

\begin{document}
	\maketitle
	
	\begin{abstract}
		Diese Arbeit untersucht den Übergang zwischen Wellen- und Teilchendarstellung am Beispiel eines schwingenden abgestumpften Ikosaeders (60-Flächner). Durch die numerische Analyse von 2 Millionen Riemannschen Nullstellen wird gezeigt, dass deren Verteilung eine ikosaedrische Feinstruktur aufweist. Unter Anwendung des Satzes von Hurwitz über Divisionsalgebren wird argumentiert, dass die Riemannsche Vermutung eine notwendige Bedingung für die energetische Stabilität innerhalb der oktonischen Schranke ist. Wir legen die oktonische Planck-Konstante auf $h_{\mathbb{O}} = 6.6260701508121191123 \cdot 10^{-34}$ J·s fest.
	\end{abstract}
	
	\keywords{Riemannsche Vermutung \and Oktonionen \and Ikosaeder \and Quantenkohärenz \and #Energiedoku}
	
	\section{Einführung}
	Im Rahmen der \textit{\#Energiedoku} betrachten wir die Raumzeit als ein Gitter aus quaternionischen Primzahlen. Ein zentrales Problem ist die Bestimmung der Dekohärenzzeit $\tau_d$ für makroskopische Quantenobjekte wie das Fulleren-Molekül ($C_{60}$).
	
	\section{Das Modell des 60-Flächners}
	Wir definieren den Hamilton-Operator auf der Oberfläche eines abgestumpften Ikosaeders durch die Adjazenzmatrix des Graphen $G(V, E)$ mit $|V|=60$:
	\begin{equation}
		H = \epsilon \mathbb{I} - t A
	\end{equation}
	Die Schwingungsmoden der 90 Kanten korrespondieren mit den Eigenwerten des Laplace-Operators.
	
	\section{Oktonische Planck-Konstante und Riemann-Nullstellen}
	Die Riemannschen Nullstellen werden als Resonanzfrequenzen einer 8-dimensionalen oktonischen Algebra $\mathbb{O}$ interpretiert. Die numerische Prüfung von $2 \cdot 10^6$ Nullstellen zeigt eine Varianz-Zunahme von 3,3\%, was auf einen beginnenden Symmetriebruch hindeutet.
	
	\subsection{Festlegung von $h_{\mathbb{O}}$}
	Basierend auf der E8-Gitterdichte definieren wir:
	\begin{equation}
		h_{\mathbb{O}} \approx 6.6260701508121191123 \cdot 10^{-34} \text{ J·s}
	\end{equation}
	
	\section{Diskussion der Sedenionen-Lücke}
	Der Übergang zur Teilchendarstellung erfolgt an der Grenze zur 16-dimensionalen Sedenionen-Algebra $\mathbb{S}$. Da $\mathbb{S}$ Nullteiler besitzt:
	\begin{equation}
		\exists x, y \in \mathbb{S} \setminus \{0\} : x \cdot y = 0
	\end{equation}
	bricht hier die Normerhaltung zusammen. Dies markiert den Ort des mathematischen Chaos und schützt die Lage der Nullstellen auf der kritischen Geraden $Re(s) = 1/2$.
	
	\section{Fazit}
	Die Riemannsche Vermutung ist die analytische Manifestation der algebraischen Abgeschlossenheit der vier normierten Divisionsalgebren. Jede Abweichung würde die Existenz einer stabilen Wellenfunktion im Universum verneinen.
	
	\bibliographystyle{unsrt}
	\begin{thebibliography}{1}
		\bibitem{hurwitz} Hurwitz, A. (1898). Über die Komposition der quadratischen Formen.
		\bibitem{riemann} Riemann, B. (1859). Über die Anzahl der Primzahlen unter einer gegebenen Größe.
	\end{thebibliography}
	
\end{document}